% 第 1 章　世界不是公式：物理学究竟在做什么（完整正文）
\chapter{世界不是公式：物理学究竟在做什么}

很多人以为，物理学是一本写满公式的书，学物理就是把这些公式背下来，再往题目里套。这本书想做一件相反的事：先把公式放一边，回到公式出现之前的那个现场——一个人盯着某个现象发呆，想不通，于是开始测量、猜测、犯错、再测量。公式是这个过程的终点，不是起点。这一章不涉及任何具体的物理定律，我们要解决一个更根本的问题：物理学究竟在做什么？它凭什么说自己“知道”？

\step{反常识开场}

请你想象一个场景，最好待会儿真的去做。甲捏着一把三十厘米直尺的上端，让尺子竖直悬空；乙把拇指和食指张开，等在尺子零刻度的两侧，不许碰到尺子。甲毫无预告地松手，乙尽快合拢手指夹住尺子，读出夹住位置的刻度。

第一次：十八厘米。第二次：十一厘米。第三次：二十三厘米。第四次：十六厘米。

奇怪的事情来了。同一个人，同一副神经和肌肉，同一把尺子，夹住的位置却一次一个样。你可能以为是乙不认真，那就让乙打起十二分精神再测十次——读数还是散开在各处。你可能会说，那就取个平均吧，平均值应该就是“乙真正的反应快慢”。可是第二天再测，平均值变了；睡个午觉再测，又变了。让班里反应最快的体育委员来测，他的读数也一样散开，只是整体偏小。

这件事看起来小得不值一提，但它戳中了物理学的软肋：我们习惯以为，世界上每个量都有一个确定的“真值”，安静地躺在那里等我们读出来，就像身高、鞋码一样。可“你的反应时间”这个量，好像根本不肯固定下来。类似的麻烦到处都是：同一块体重秤，早上和晚上能差出一公斤；同一段上班的路，导航软件每天预计的到达时间都不一样；百米决赛的电子计时精确到千分之一秒，可要是让同一名运动员再跑一次，成绩从来不会一模一样；手机地图上代表你位置的蓝点，站在原地不动时也会自己慢慢漂移，每次都给出稍有不同的坐标。

世界给我们的从来不是“那个数”，而是一堆彼此不太一样的数。在这样乱糟糟的数据面前，“知道”到底是什么意思？物理学家凭什么敢下结论？这就是本章要回答的问题。

\step{先猜一猜}

在往下读之前，请把你的直觉预测写下来，越明确越好。后面我们会逐一回来对账。

第一，如果你认真测十次反应时间，最快和最慢大约会差多少？你的直觉是“差不到百分之几”，还是“可能差出一倍”？

第二，如果甲在松手前喊一声“预备”，乙的成绩会变好吗？读数的散布会变大还是变小？

第三，换左手测和换右手测，结果应该一样吗？

第四，也是最关键的一猜：如果测一百次、一千次，再取平均，能不能得到“真正的反应时间”，精确到小数点后三位？

大多数人凭直觉回答第四问时会说：能，只要测得足够多次，平均值就会无限逼近真值。这个直觉有一半对，有一半会栽跟头。栽跟头的地方，正是本章最重要的一课。请把你的四个答案记在纸角上——物理学里最值钱的不是正确答案，而是一个说清楚了、可以被判决的猜测。

\step{动手实验}

\begin{experiment}[尺子反应时间]
\textbf{材料}：一把三十厘米直尺（塑料的就行）、纸和笔、手机计算器。

\textbf{步骤}：
\begin{enumerate}[leftmargin=2em]
  \item 请家人或同学捏住直尺上端，让尺子竖直下垂，零刻度在下方。
  \item 你的拇指和食指张开约三厘米，虚拢在零刻度两侧，\textbf{不要碰到尺子}。
  \item 对方在不预告的时刻松手，你尽快夹住，读出夹住处的刻度，记下来。
  \item 重复十次，每次都把尺子重新提回原来的位置。
  \item 算出十次读数的平均值，再找出最大值和最小值，记下它们的差（这个差叫\textbf{极差}）。
\end{enumerate}

\textbf{参考数据}：笔者的一次实测是 \(12, 15, 16, 17, 18, 18, 19, 21, 23, 27\)（单位：厘米），平均值 \(18.6\) 厘米，极差 \(15\) 厘米。你的数据几乎肯定不一样，但如果最快和最慢相差不到两三厘米，反而值得怀疑是不是每次都有预告。

\textbf{变体}：许多手机应用商店里有“反应时间测试”小程序，用触屏代替尺子，可以直接读出毫秒数。试一试，比较两种做法给出的散布哪个更大，想想为什么。
\end{experiment}

\begin{center}
\begin{tikzpicture}[scale=0.9]
  \draw[thick] (0,0) rectangle (0.9,5.4);
  \foreach \y in {0,0.9,...,5.4} \draw (0,\y) -- (0.22,\y);
  \node[left] at (0,0) {\small \(0\)};
  \node[left] at (0,5.4) {\small \(30\)};
  \node[rotate=90] at (0.45,2.7) {\small 直尺};
  \draw[thick] (-0.35,5.05) -- (0,5.05);
  \draw[thick] (1.25,5.05) -- (0.9,5.05);
  \node[above] at (0.45,5.5) {\small 甲捏住上端，随机松手};
  \draw[thick] (-0.6,0.15) -- (0,0.15);
  \draw[thick] (1.5,0.15) -- (0.9,0.15);
  \node[below] at (0.45,-0.15) {\small 乙的手指虚拢在零刻度两侧};
  \draw[dashed] (0,3.3) -- (1.75,3.3);
  \draw[<->,>=Stealth] (1.75,0) -- (1.75,3.3) node[midway,right]{\small 读数 \(h\)};
\end{tikzpicture}
\end{center}

这个实验没有任何危险的环节，却包含了全部物理研究的骨架：有一个现象（尺子被夹住的位置），有一套操作约定（怎么捏、怎么等、怎么读数），有一串数据（十次读数），还有一个诱惑——想从数据里挤出一个“结论”。接下来我们看看，直觉在这串数据面前会遇到什么麻烦。

\step{旧模型遇到麻烦}

面对十次不同的读数，直觉会搬出两个旧模型来救场。我们一个一个看它们怎么倒下。

\textbf{旧模型一：“每个量都有一个真值，测量就是把它读出来，读不准就是出了错。”} 这个模型暗示：十次读数里至多有一个是对的，其余九次都是“错误”。可是你看看数据的形状——多数读数挤在中间，偏离很远的很少，这不是“九次犯错”该有的样子，倒像是有某种规律在支配散布。更麻烦的是，散布不会因为你更认真而消失。换上带电子计时和光电传感器的精密装置，让同一个人反复做同样的反应测试，读数的最后一位仍然会跳动。如果散布只是“粗心”，它应该能被认真消灭；事实是它只能被压缩，不能被消灭。这个旧模型还欠下一笔债：它说不出“多读数几次取平均”为什么有意义，也说不出该报到小数点后第几位。一个连操作指南都给不出的模型，留着何用。

\textbf{旧模型二：“学物理就是背公式，公式就是世界本身。”} 这个模型更深，也更顽固。给它致命一击的是一个历史事实：托勒密的地心说认为太阳和行星都绕地球转，行星还沿着一套套嵌套的小圆（本轮）运行。这个模型今天我们认定是错的，可是只要耐心地调整那些小圆的大小和转速，它预言行星在天空中的位置，精度足以指导一千四百年间的历法、航海和占星——在很多实际用途上，它“非常好用”。一个错误的模型可以在有限的范围内交出漂亮的答卷。反过来说，“好用”不能证明一个模型就是世界的本来面目。公式不是世界，公式是我们为世界画的像；画像可以神似，可以失真，也可以在这个角度神似、换个角度失真。

还有一个日常版本的老观念值得顺手拆掉：“测得越精细越好。”直觉说，报数字报的位数越多，测量水平越高。反例来了：用普通米尺量身高，尺子上最小的格子是毫米，有人却报出“一百七十二点三四一八厘米”——小数点后第三、四位根本不是量出来的，是算平均值时计算器吐出来的。更典型的反例是一台没调零的体重秤：指针空载时指在两公斤处，你站上去测一万次再取平均，平均值依然稳稳地偏出两公斤。直觉说“测一万次总该准了吧”，错了——多次测量能压缩的是那种忽大忽小的起伏，对这种每次都朝同一个方向偏的毛病，平均一万次也无能为力。这两种毛病必须分开对待，这正是下一步要做的事。

\step{发明新概念}

旧模型倒下的地方，就是新概念出生的地方。物理学用了一组分工明确的词，把“知道”这件事拆成了几道工序。请先记住这五个词的区别，全书都要靠它们吃饭。

\begin{mainline}
\textbf{事实}：世界实际发生的事，以及可重复、可核查的观察记录。“松开手的苹果向下落”“乙十次夹尺的读数是这些数”——这是事实层。事实不依赖任何理论，但记录事实需要约定（用什么仪器、从哪读数）。

\textbf{测量}：用仪器和约定，把现象翻译成数。测量总是有“分辨率”的：米尺读到毫米，千分尺读到百分之一毫米，没有仪器能给出无限多位有效数字。

\textbf{模型}：我们发明的、关于世界如何运作的简化故事，通常配着数学。模型主动扔掉大量细节，只留下它认为要紧的关系。

\textbf{定律}：模型中被大量实验反复检验、在明确范围内从未失手的关系，例如“在地球表面附近，一切物体下落的快慢与轻重无关”。定律依然属于模型层，不属于事实层，只是模型里成绩最好的那一部分。

\textbf{解释}：把一件事实放进某个模型里讲通的过程。“苹果下落是因为地球引力拉它”是解释，不是事实本身。同一个事实可以有几种解释互相竞争，靠实验的新预言来裁判。
\end{mainline}

拿尺子实验把这五道工序各过一遍。事实：纸上那十个读数，以及“读数总是散开”这个观察。测量：从零刻度读到夹住处，约定估读到半厘米。模型：“反应时间有一个典型水平，每次测量在这个水平附近随机起伏”——这句话本身就是一个模型，它扔掉了你那天早饭吃了什么、窗外有没有车按喇叭等一切细节。定律层本章还谈不上，但到了第 5 章，“下落距离与时间的平方成正比”就会成为你能检验的定律。解释：读数偏大的一次，可以解释为“那次注意力飘走了”；注意，这个解释对不对，要靠新实验（比如先给预备口令再看读数是否整体变小）来裁判，而不是靠它听起来合理。

回到苹果。“物体会下落”是可以反复核查的事实；“为什么下落”则需要模型——你可以说“地球用引力拉它”（牛顿的模型），也可以说“地球弯折了周围的时空，苹果沿着弯曲时空里最直的路径走”（爱因斯坦的模型）。两者都解释了下落，后者还能解释前者解释不了的细节，比如水星轨道的微小偏差。哪个是“真相”？物理学家会换个说法：哪个模型的预言与实验对账对得更准、范围更大，哪个就暂时领先。模型没有终身成就奖，只有逐年的成绩单。

接下来是三个全书通用的理想化概念，它们从下一章起会天天露面。

\textbf{质点}：研究地球绕太阳公转时，把直径一万两千多公里的地球当成一个没有大小的点。荒谬吗？可是公转轨道的半径约一亿五千万公里，地球的“胖瘦”只占行程的万分之一量级，扔掉它，计算立刻从不可能变成几行算式，而预言的轨道位置几乎不变。

\textbf{光滑平面}：真实桌面总有摩擦，但研究“物体不受推挤时会怎样”这类问题时，摩擦会把最核心的规律掩盖起来。先想象一个摩擦力为零的平面，把规律看清楚了，再把摩擦作为一项修正加回来。第 6 章讲牛顿第一定律时你会看到，没有这个理想化，人类可能至今还在相信“物体需要力才能维持运动”。

\textbf{不可伸长的绳}：真实绳子都会被拉长一点，但研究摆、滑轮时，忽略这点伸长，问题立刻清爽。等到研究蹦极——绳子的伸长恰恰是保命的全部关键——这个理想化就必须第一个被扔掉。

这些理想化为什么有用？这里用本书第一个正式比喻。\textbf{模型像地图}：它像地图的地方在于，一幅把每条街道、每棵树都按一比一复制的地图毫无用处，地图的价值恰恰来自舍弃——舍弃你不关心的细节，突出你关心的关系；它不像地图的地方在于，地图只负责描摹已经存在的地方，而物理模型必须走得更远：它要对还没做过的实验、还没测过的情形做出预言，接受判决。一幅永远不会被新测量打分的模型，无论多漂亮，都不是物理学。

最后给测量中的两种“麻烦”正式命名，刚才那台没调零的秤你已经见过它们了。一种叫\textbf{随机起伏}：每次测量忽大忽小，方向没准，取平均可以把它压下去；一种叫\textbf{系统偏差}：每次都朝同一个方向偏，比如秤没调零、尺子受热变长、你读数时总习惯从斜上方看——取平均对它无效，只能靠改进仪器和操作方法来消除。测量结果的品质，用\textbf{不确定度}来刻画：它不是“我错了”的道歉，而是“真值大概率落在哪个区间”的量化声明。与不确定度配套的是\textbf{有效数字}的规矩：测量结果报到哪一位，要看仪器分辨率和数据散布允许到哪一位；超出这个范围的位数，是自欺欺人的装修。一个可操作的经验法则是：结果的末位应与不确定度同数量级。十次夹尺读数的极差是十五厘米，平均值报“十八点六厘米”已经勉强，报“十八点六二四厘米”就是给计算器位数镀金；反过来，精密天平测到零点零零一克，你却只报“约两克”，同样是浪费——把仪器好不容易挣来的信息扔掉了。

\step{一句话模型}

\begin{oneline}
物理学不是背公式，而是一套流程：盯住现象，发明简化模型，从模型推出可检验的预言，再让实验打分；模型没有“真”或“假”的证书，只有“在哪个范围内、准到什么程度”的成绩单。
\end{oneline}

\step{数学翻译器}

直觉说“取个平均看看”，数学问：平均是什么？凭什么它代表十次测量？我们把第一句话翻译成符号。设十次读数为 \(x_1, x_2, \ldots, x_n\)（这里 \(n=10\)），平均值定义为
\[
\bar{x}=\frac{x_1+x_2+\cdots+x_n}{n}\,.
\]

用本章的规矩，每条公式都要回答四个问题。\textbf{它描述什么}：一组测量值的“中心位置”。\textbf{为什么这样写}：加起来再平分，是唯一对每个读数一视同仁、且满足“整体抬高多少，平均值就抬高多少”的算法；任何偏心某个读数的算法都没资格代表全体。\textbf{何时成立}：当你想概括的是“同一条件下的反复测量”，并且数据里没有混入别的东西（比如中途换了个人来测）。\textbf{何时失效}：当数据本身来自两个不同的分布——比如你测了五次、换反应更快的同桌又测了五次，混合平均出的数既不代表你也不代表他；当存在系统偏差时，平均值再精确也忠实地把偏差继承下来。

再用特殊情况检查它，这是本书的固定动作。\textbf{全相同}：十次读数都是 \(18\)，平均值就是 \(18\)，合理。\textbf{整体平移}：每个读数都加 \(2\)（尺子每次都往下多放了两厘米），平均值恰好加 \(2\)，说明它如实跟随整体变化。\textbf{单个离群}：某一次读数手滑变成 \(50\) 厘米，平均值会被明显拉高——这提醒你平均值怕离群值，出现异常数据要先查操作，而不是闷头照算。\textbf{\(n=1\)}：只测一次时平均值就是这一次，公式退化得自然，但也暴露了一个事实——单次测量完全无法告诉你散布有多大。

把笔者那组数据画成“读数落在各区间里的次数”，形状一目了然：中间高、两边低，像一座小山头。

\begin{center}
\begin{tikzpicture}[xscale=1.1,yscale=0.5]
  \draw[->,>=Stealth] (0,0) -- (5.4,0) node[right]{\small 读数 \(h\)（厘米）};
  \draw[->,>=Stealth] (0,0) -- (0,7) node[above]{\small 次数};
  \draw[fill=black!20] (0.2,0) rectangle (1.2,1);
  \draw[fill=black!20] (1.2,0) rectangle (2.2,6);
  \draw[fill=black!20] (2.2,0) rectangle (3.2,2);
  \draw[fill=black!20] (3.2,0) rectangle (4.2,1);
  \node[below] at (0.7,0) {\small 10--15};
  \node[below] at (1.7,0) {\small 15--20};
  \node[below] at (2.7,0) {\small 20--25};
  \node[below] at (3.7,0) {\small 25--30};
  \foreach \y in {1,...,6} \node[left] at (0,\y) {\scriptsize \y};
\end{tikzpicture}
\end{center}

这座“小山”就是物理学家说的\textbf{分布}。它才是测量真正的主角：平均值只是山的重心位置，而山的宽度、形状同样携带信息。第二册讲热学时你会再见到它——亿万分子的速率也排成这样一座山；第三册讲量子时，它还会以更深刻的方式回来。

散布怎么翻译？主线层先用极差：最大值减最小值，笔者那组数据是 \(27-12=15\) 厘米。极差简单，但太依赖两个极端值。更稳的度量放在数学桥层。

\begin{mathbridge}
把每个读数与平均值的偏差平方、平均、再开方，得到\textbf{标准偏差}：
\[
u=\sqrt{\frac{1}{n-1}\sum_{i=1}^{n}\left(x_i-\bar{x}\right)^2}\,.
\]
\textbf{描述什么}：读数围绕平均值的典型起伏幅度。\textbf{为什么这样写}：偏差有正有负，直接相加会正负抵消，平方后全变正，大起大落被放大、小起小落被压低；开方让单位回到原来的厘米。\textbf{何时成立}：单次测量的起伏彼此独立、没有明显趋势。\textbf{何时失效}：数据若有趋势（越测越顺手，读数单调变小），标准偏差会把“进步”误算成“起伏”。

检查特殊情况：全部读数相同，每个偏差都是零，\(u=0\)——散布为零，合理；所有读数同加 \(2\)，偏差不变，\(u\) 不变——散布不受整体平移影响，合理；单位从厘米换成毫米，每个偏差乘十，\(u\) 也乘十——它像一把真正的尺子，随单位缩放。
\end{mathbridge}

\begin{challenge}
公式里除以的是 \(n-1\) 而不是 \(n\)，这是贝塞尔修正。直觉是：\(\bar{x}\) 本身就是从这 \(n\) 个数里算出来的，已经“借用”了数据的一份信息，\(n\) 个偏差里真正独立的只剩 \(n-1\) 个——事实上 \(n\) 个偏差之和恒等于零，知道前 \(n-1\) 个，最后一个就被定死了。除以 \(n\) 会系统性地把散布估小。严格的证明需要一点概率论，这里只要记住：从数据里借用过的东西，算散布时要还回去。
\end{challenge}

现在兑现一个真实数据的估算。尺子读数怎么换算成反应时间？换算需要自由下落的规律，那是第 5 章的任务，这里先只引用结论：从静止松手，尺子下落 \(h\) 所用的时间是
\[
t=\sqrt{\frac{2h}{g}}\,,\qquad g\approx 9.8\ \mathrm{m/s^2}\,.
\]
代入 \(h=18.6\) 厘米 \(=0.186\) 米，得 \(t\approx\sqrt{2\times0.186/9.8}\approx 0.19\) 秒。这和电子设备直接测出的人的反应时间（约 \(0.15\) 到 \(0.25\) 秒）吻合。检查特殊情况：\(h=0\) 时 \(t=0\)，手指本来就在零刻度处，零时间，合理；\(h\) 变成四倍，\(t\) 只变成两倍——\textbf{时间不与下落距离成正比}。这一点直觉经常出错：同桌夹尺读数是你的两倍，他的反应时间并不是你的两倍，只是约 \(\sqrt{2}\approx1.4\) 倍。拿尺子读数直接比较反应快慢，会冤枉人。

这个估算马上有现实用处。汽车时速一百公里，约每秒二十七点八米；驾驶员从看见险情到踩下刹车大约要零点五秒（含判断时间，比夹尺子慢得多），这期间车已经冲出约十四米——这就是交通法规里“安全车距”的物理来源。自动驾驶系统把这段反应压缩到零点一秒以内，省下的不只是零点四秒，而是每小时一百公里速度下的十多米生命空间。工程设计里的“反应时间预算”，正是从你我夹尺子的那个分布里长出来的。

\step{模型的边界}

现在给本章的概念一一划出边界，这是比发明概念更重要的自律。

\textbf{理想化的失效清单}。质点模型能算地球公转，但研究地球自转、昼夜、潮汐时，地球的大小和形状就是主角，质点立即破产；光滑平面能让运动规律现形，但冰壶运动员拼命刷冰，正是在调制摩擦，把摩擦抹掉你就永远看不懂冰壶；不可伸长的绳对秋千是好朋友，对蹦极绳是谋杀——蹦极的全部安全性恰恰来自绳子的伸长把冲击力摊到更长的时间里。理想化不是谎言，是带说明书的工具：说明书上写着“本工具忽略了什么”，用之前必须读。

\textbf{错误模型的合法使用区}。托勒密体系错了，但指导一千四百年的历法足够好；“地面是平的”错了，但盖房子、修操场时把大地当平面，误差小到可以忽略；把空气阻力当不存在错了，但分析从一米高处落下的书本时，忽略它没什么后果。一个模型在什么范围内准到什么程度，是模型知识的一部分。物理学的成熟，不在于找到终极正确的模型，而在于给每个模型配上清晰的适用范围——超出范围的外推，是工程事故和科学事故共同的温床。

天气预报是这条原则最公开的演示。气象中心把大气切成几千万个小格子，用流体、热量和水汽交换的方程推算格子之间的相互作用，让超级计算机把这个巨大的模型往前演化几天，再翻译成你手机上的降水概率。这个模型显然不是真实的大气——真实大气没有格子——但只要格子的加密、方程的改进被观测数据持续检验，它就在自己的适用范围内越来越好用；同时每个人都知道，七天以后的预报要打折看，这正是公众在不知不觉中学会了阅读“模型的适用范围”。

\textbf{统计量的误用区}。平均值代表“典型水平”，不代表“每一次的值”——你的平均反应时间是零点二秒，不等于今晚疲劳驾驶时也是零点二秒；散布大本身就是重要信息，只报平均值不报散布，等于隐瞒病情。有效数字报得再长，也长不过仪器分辨率和数据散布允许的范围。

\textbf{一个必须立刻澄清的混淆}。本章讲的测量不确定度，是经典意义上的：它与仪器分辨率、环境扰动、人的状态有关，原则上可以通过改进技术不断压缩。第三册讲量子力学时会遇到著名的“不确定性原理”——名字相似，实质完全不同：那条原理说的是，某些成对的物理量（例如位置和动量）天然不能同时具有任意精确的值，这不是仪器粗糙造成的，再完美的仪器也无法绕过，它是自然本身的结构。把两者混为一谈，是科普中最常见的错误之一，本书从现在起就严格区分：本章的是“测得准不准”的问题，那是“存不存在一个同时确定的值”的问题。

\textbf{系统偏差的顽固区}。随机起伏能被多次测量压缩，系统偏差不能。没调零的秤测一万次还是偏两公斤；用受热膨胀的尺子量布，全年都在偏。消除系统偏差没有统计学捷径，只有一条笨路：换不同的仪器、不同的方法、不同的实验室交叉对账。科学史上许多著名乌龙，都是败在没人怀疑过的系统偏差上。

\step{回头解释开场现象}

现在，带着新概念回到那把直尺面前。

开场时我们以为的麻烦——“同一个人的反应时间为什么每次不同”——在新语言里不再是麻烦，而是再正常不过的情形。“你的反应时间”从来不是一个单一的数，它是一个分布：有一个中心（平均值，约零点二秒），有一个宽度（标准偏差，通常几十毫秒），形状中间高、两边低。每次测量是从这个分布里抽一次样，读数散开不是测量出了错，而是被测的东西本来就有宽度。神经信号传导的微小快慢、注意力的漂移、手指启动的差异，都真实地存在于每一次反应里——仪器再精密，也只能如实报告这份宽度。

先猜一猜里的四问，现在可以对账了。第一，十次读数差出一倍很正常，因为散布来自人本身，不是尺子的毛病。第二，“预备”口令会缩短平均反应时间、收窄散布，因为注意力被集中了——这是模型给出的可检验预言，你回家就能验证。第三，惯用手通常略快，差别约在散布的量级，需要成对比较多次才能看出来。第四，也就是直觉栽跟头的那个：测一千次取平均，平均值的确会稳定下来，但稳定的只是平均值，单次测量的散布一点都不会缩小；更要命的是，如果夹尺方法本身有系统偏差（比如你总从斜上方读数），平均值会稳定地停在错误的位置上。无限多次测量能把随机起伏压到任意小，却拿系统偏差毫无办法——这就是为什么物理学不能只靠“多测”，还要靠“换着法子测”。

开场悬念至此解开：测量结果每次不同，不是因为物理学无能，而是因为“反应时间”这个概念本身就是对一堆起伏的概括。物理学不假装世界是钉子一样钉死的数，它发明分布、平均值、不确定度这些概念，老老实实地刻画世界的起伏，并从中提取稳定可靠的规律。顺带说一句，这套语言不只属于实验室：体检报告上的参考区间、民意调查的“误差范围”、药品说明书上的有效率，全都是“分布加不确定度”的写法。学会这一章，你已经能拆穿生活里一大半打着数字旗号的虚张声势。

\step{脑洞实验与挑战题}

\begin{thought}[把地球装进一个点]
本章说，研究地球公转时可以把直径约一万两千七百四十二公里的地球当成一个质点。这个近似到底有多“好”？日地距离约一亿五千万公里，地球直径只占它的大约万分之一——也就是说，把一个庞然大物捏成没有大小的点，给公转轨道带来的位置不确定度只在万分之一的量级，远小于古代天文观测的精度，甚至连十九世纪的观测都挑不出毛病。

\begin{center}
\begin{tikzpicture}
  \draw[fill=yellow!50] (0,0) circle (0.45);
  \node[below=10pt] at (0,-0.45) {\small 太阳};
  \draw[dashed] (0,0) circle (3);
  \draw[<->,>=Stealth] (0,0) -- (2.12,2.12) node[midway,above left=-2pt]{\small 约 \(1.5\) 亿公里};
  \fill (3,0) circle (0.05);
  \node[right] at (3.05,0) {\small 地球（当成一个点）};
\end{tikzpicture}
\end{center}

现在把镜头推向极端：假如未来天文学家要精确预言一万年后的日食——阳光被月亮挡住时，影锥扫过地球表面，落在哪个城市、哪条街道。这时候把地球当质点会闹出灾难：影子在地表的轨迹宽度只有几百公里，而“地球是个点”等于宣布地球的位置有一万两千公里的模糊。同一个理想化，在公转问题里是天才，在日食预言里是废物。模型的价值从来不写在模型身上，而写在你拿它去回答的那个问题里。
\end{thought}

\begin{puzzles}
\item 小明用普通直尺测十次反应时间，小红用带光电传感的精密电子装置也测十次。小红的十次读数，散布一定比小明小吗？

\textbf{思路提示}：先分清散布的两个来源——仪器分辨率带来的，和被测对象本身的起伏带来的。反应时间的散布大头在哪里？换精密仪器能压缩的是哪一部分？
\item 某实验室报告：“我们测得这种材料的强度是 \(137.4281\) 兆帕，重复测量十次的极差是 \(9\) 兆帕。”这份报告里有一个自相矛盾的地方，找一找。

\textbf{思路提示}：平均值的有效数字位数，有没有资格超过数据散布允许的范围？极差为 \(9\) 兆帕时，小数点后第四位在说什么？
\item 有人据此嘲讽物理学：“既然你们承认每次测量都不一样，那物理学什么也不能真正确定，科学不过是一种信仰。”请用本章的概念反驳，不超过五句话。

\textbf{思路提示}：散布本身有没有规律？不确定度是不是一个可以量化、可以压缩、可以交叉验证的量？“不能确定到任意多位数”和“什么都不能确定”之间隔着什么？
\item 给你一把最小刻度为毫米的直尺，要求估出一张纸的厚度，并说出你的不确定度。直尺直接测一张纸毫无办法，怎么办？这种办法对系统偏差有效吗？

\textbf{思路提示}：一张纸测不了，五百张叠起来呢？把总厚度除以张数，相当于把仪器分辨率“摊薄”了多少倍？再想想：如果这叠纸里混着一张明显更厚的，或者尺子受热变长了，平均能救你吗？
\end{puzzles}

本章结束时，请把贯穿全书的三个追问抄在笔记本第一页，它们将从下一章开始反复出现：\textbf{系统此刻是什么状态？什么规律决定状态如何变化？我们怎样通过测量知道它？}反应时间的例子已经提前给出了答案的雏形：状态是“一个分布”，规律是“分布的中心和宽度如何随条件改变”，测量是“抽样、取平均、评估不确定度、排查系统偏差”。等到第三册我们讨论量子世界的概率、第二册讨论亿万分子的统计规律时，你会再次回到这把直尺面前，发现第一章的麻烦早已预告了整部物理学的语法：世界从不直接交出答案，它只交出数据，而物理学是解读数据的纪律。
